% s3: the note of one turn (strong coupling). ASCII only.
\section*{3. The note of one turn}

Everything physical about a gauge field is carried by its turns: carry a frame around a loop and see how far it comes
back rotated. In the purely electric strong-coupling limit the first excitation is the
smallest closed flux loop. Wilson's plaquette term $1-\tfrac12\mathrm{tr}\,U_p=1-\cos\alpha$
is half the \emph{squared} chord distance between the unit quaternions $U_p$ and $1$; in \emph{Three Ways Across a Turn}~\cite{ThreeWays}
I compare it with the arc and the tangent. The cleanest view is the Hamiltonian lattice of Kogut and
Susskind~\cite{KogutSusskind75}: continuous time, a cubic spatial lattice, and at strong coupling an energy that is
\emph{electric}, $g^2/(2a)$ times the quadratic Casimir $C_2$ on each link. The Casimir is the Laplacian of the group,
and $SU(2)$ is a three-sphere, a closed curved shape with a lowest nonzero note: $C_2(j)=j(j+1)$, smallest at
$C_2(\tfrac12)=\tfrac34$, which exists only because $SU(2)$ double covers $SO(3)$ (on $SO(3)$ alone it would be $2$).
Gauss's law allows electric flux only in closed loops, and the shortest is one plaquette.

\begin{claim}{PROVED}{The electric gap is one turn}\label{b:electric}
On the cubic lattice (infinite, or periodic with side $L\ge4$) with no external charges, $H_E=(g^2/2a)\sum_\ell C_2(\ell)$
has a unique ground state of energy $0$ and first excited level exactly $4\cdot\tfrac{g^2}{2a}\cdot\tfrac34=3g^2/(2a)$:
one plaquette loop in spin $\tfrac12$.
\evidence{Peter--Weyl splits gauge-invariant states by link spins; excited links must form cycles of at least four edges
\cite{KogutSusskind75,HelikeReport}.}
\lines{4}
\end{claim}

The magnetic term, of order $1/(ag^2)$, mixes in other configurations, so $3g^2/(2a)$ is the leading term of the gap. Two
lattices must be kept apart here (I once mixed them; the register records the correction~\cite{ClosedSet}). On the
\emph{Euclidean} lattice, with discrete time, a strong-coupling gap is a theorem; on a fixed finite
\emph{Hamiltonian} lattice the leading value follows by perturbation.

\begin{claim}{PROVED (CITED)}{A gap at strong coupling}\label{b:strong}
For small $\beta$ the cluster expansion converges and truncated correlations of local gauge-invariant observables decay as
$e^{-m(\beta)\cdot\text{distance}}$, $m(\beta)>0$; by reflection positivity this is a gap of the transfer matrix.
\evidence{Osterwalder and Seiler~\cite{OsterwalderSeiler78}; by another road, Shen, Zhu and Zhu~\cite{ShenZhuZhu23}.}
\lines{3}
\end{claim}

\begin{claim}{PROVED}{The leading Hamiltonian gap (by deduction)}\label{b:kato}
On a finite Hamiltonian lattice the magnetic term is a bounded perturbation of $H_E$, so the full Kogut--Susskind gap is
$3g^2/(2a)+O_L(1/(ag^2))$ as $g\to\infty$. The constant here may depend on lattice size $L$.
\evidence{Kato's analytic perturbation theory~\cite{Kato}. A volume-uniform Hamiltonian statement needs a separate
uniform argument; the Euclidean result in Box~\ref{b:strong} is not by itself that argument.}
\lines{5}
\end{claim}

In two dimensions the whole story is exact at every coupling: in a tree gauge the plaquettes are independent and every
loop factorises~\cite{Migdal75,GrossWitten80}.

\begin{claim}{PROVED}{Two dimensions: each loop's area law is its Casimir}\label{b:2d}
For $SU(2)$ on a planar lattice, a simple loop around $|A|$ plaquettes has $\langle\chi_j\rangle/d_j=(I_{2j+1}(\beta)/
I_1(\beta))^{|A|}$, so $\sigma_j=2C_2(j)/\beta+O(\beta^{-2})$ in lattice area units and $\sigma_1/\sigma_{1/2}\to8/3$. A centre vortex (links
crossed by a sheet multiplied by $-1$) multiplies a spin-$j$ loop by $(-1)^{2jn}$, $n$ its winding around the vortex:
half-integer loops feel it, integer loops do not.
\evidence{Character factorisation~\cite{Migdal75,GrossWitten80}; $\chi_j(-U)=(-1)^{2j}\chi_j(U)$, the centre-vortex
picture~\cite{tHooft78,Greensite03}. Casimir scaling is seen in four dimensions~\cite{Bali00}.}
\lines{6}
\end{claim}

I checked the two-dimensional loop identities in GIGI, my database engine, with the lattice connection stored as a bundle and every loop read
only by point queries: the loops match the character formula within $2\sigma$, a random gauge rewrite changes no loop by
more than $1.23\times10^{-15}$, and a planted vortex flips the fundamental loop exactly and leaves the adjoint
alone~\cite{HelikeReport}.

Strong coupling gives the note but not the continuum, and it is never an end of the family of rooms, since every
room's continuum limit sends the bare coupling to zero (\S7). Two dimensions has an exact area law but no propagating
gluons. Near the identity the Wilson plaquette cost is quadratic in angle, with a coupling-dependent coefficient;
that local cost alone does not determine a quantum gap. So in
four dimensions the note has to be found again, carried by a shape, in a room small enough that the coupling is weak.
